\documentclass[../../../include/open-logic-section]{subfiles}

\begin{document}

% Part: first-order-logic
% Chapter: beyond
% Section: many-sorted-logic

\olfileid{fol}{byd}{msl}

\olsection{多类逻辑}

在一阶逻辑中，变元和量词仅在单一论域上遍历。但拥有多个（互不相交的）论域通常很有用：例如，我们可能需要一个数的论域、一个几何对象的论域、一个从数到数的函数的论域、一个阿贝尔群的论域，等等。

多类逻辑恰好提供了这样的框架。首先给出一个“类”的列表——对象的“类”表明了它应当属于哪个“论域”。然后，为每个类引入变元和量词，并且（通常）为每个类引入一个等词谓词。函数和关系也根据它们所取自变元的类而被“赋予类型”。除此之外，我们保留一阶逻辑的通常规则，只是为每个类各重复一版量词规则。

例如，为了研究国际关系，我们可能会选择一种包含两类对象的语言：法国公民和德国公民。我们可能会有一个一元关系“喝葡萄酒”，适用于第一类对象；另一个一元关系“吃香肠”，适用于第二类对象；以及一个二元关系“结成跨国夫妻”，它取两个自变元，其中第一个自变元属于第一类，第二个自变元属于第二类。如果我们让变元 $a$、$b$、$c$ 遍历法国公民，让变元 $x$、$y$、$z$ 遍历德国公民，那么
\[
\lforall[a][\lforall x][(\Atom{\Obj{MarriedTo}}{a,x} \lif
(\Atom{\Obj{DrinksWine}}{a} \lor \lnot \Atom{\Obj{EatsWurst}}{x})]]
\]
断言：若任何法国人与某个德国人结婚，则要么该法国人喝葡萄酒，要么该德国人不吃香肠。

多类逻辑可以以一种自然的方式嵌入到一阶逻辑中，方法是将所有多类论域中的对象合并成一个单一的一阶论域，使用一元谓词符号来记录对象的类，并对量词进行相对化。例如，与上述例子对应的一阶语言将拥有一元谓词符号“$\Obj{German}$”和“$\Obj{French}$”，以及前面描述过的其他关系，但原有的类要求被去除。一个带类量词 $\lforall[x][!A]$（其中 $x$ 是德国类的变元）被翻译为
\[
\lforall[x][(\Atom{\Obj{German}}{x} \lif !A)].
\]
我们需要添加一些公理来确保这些类是互不相交的——例如，$\lforall[x][\lnot (\Atom{\Obj{German}}{x} \land
  \Atom{\Obj{French}}{x})]$——以及一些公理来保证“喝葡萄酒”只对满足谓词 $\Atom{\Obj{French}}{x}$ 的对象成立，等等。有了这些约定和公理，就不难证明多类语句可以翻译为一阶语句，且多类推导可以翻译为一阶推导。此外，多类结构也可以“翻译”成相应的一阶结构，反之亦然，由此我们也得到了多类逻辑的完备性定理。

\end{document}